\chapter{Encoding a Problem in Qubits}\label{ch:encoding}

\begin{etymology}
\textbf{Encode} comes from \emph{code}, from the Latin \emph{codex}, ``tree
trunk'', then the wooden tablets used for writing, then a book, and finally a
book of laws. To encode is to choose a set of symbols and rules to write something
down.

\textbf{Register} comes from \emph{regesta}, ``things written down'', from
\emph{regerere}, ``to record''. The place where information is kept while you
work with it.

\textbf{Basis} is Greek for ``step'' or ``what you stand on''. The set of
references you describe everything else with.
\end{etymology}

\begin{quote}
To encode a problem in qubits is to decide what the basis states of a register
stand for.
\end{quote}

In practice that means choosing a correspondence between configurations of the
problem and strings of bits. A register of $n$ qubits has $2^n$ basis states, so
encoding is deciding how the variables, candidates or answers of the problem are
translated into those states. This chapter shows how that works in general, how
Shor did it for factoring, and what each choice costs.

\section{The computational basis is an alphabet}

The basis states of a register are all the strings of $n$ bits. One qubit has
$\ket0$ and $\ket1$. Two qubits have $\ket{00}, \ket{01}, \ket{10}, \ket{11}$. Three
have eight, from $\ket{000}$ to $\ket{111}$.

Imagine the eight states of three qubits are eight people. The first bit says
whether a person wears a hat, the second whether they wear glasses, the third
whether their hand is raised. Then $\ket{010}$ and $\ket{111}$ are not mysterious;
they are labels for two profiles. Each basis state is one configuration of the
problem.

A classical computer that wants to find the people matching some rule checks them
one by one. A quantum computer can put the register in a superposition of all
eight and apply the rule to all of them in a single operation: ``everyone who
matches, flip your sign''. The operation acts on every component of the
superposition at once.

That does not mean you can read the eight answers. A measurement returns one basis
state. Getting the full record of every person is a question a quantum computer
does not answer any better than a classical one. What it can do is interfere the
signs so that the measurement is likely to return a person who matches. With $n$
features there are $N = 2^n$ profiles. Checking them one by one costs $O(N)$; with
Grover's interference it costs $O(\sqrt N)$.

\section{Registers}

A register is a group of qubits that together represent a number in binary. With
three qubits, $\ket{101}$ means $1\cdot4 + 0\cdot2 + 1\cdot1 = 5$. So far this is
exactly what a classical register does. The difference is that a Hadamard on each
of the three qubits puts the register in an even superposition of all eight
values, $0$ to $7$, at once. Three physical qubits, a space of eight dimensions.
To hold the numbers from $0$ to $N - 1$ you need $\lceil \log_2 N \rceil$ qubits:
four for $N = 15$, since $2^3 = 8 < 15 \leq 16 = 2^4$.

\section{How Shor encodes factoring}

Chapter~\ref{ch:factoring} reduced factoring to finding the period $r$ of $f(a) = x^a \bmod N$. Shor
needed two things in the quantum state at the same time: the inputs $a$ being
tried, and the value of $f$ at each of them. So he used two registers.

\begin{figure}[htbp]
\centering
\begin{tikzpicture}[node distance=7mm]
  \node[hf/reader, minimum width=26mm, minimum height=10mm] (r1) {register 1\\{\footnotesize $t$ qubits, holds $a$}};
  \node[hf/source, minimum width=26mm, minimum height=10mm, below=6mm of r1] (r2) {register 2\\{\footnotesize holds $x^a \bmod N$}};
  \node[hf/bus, minimum width=12mm, minimum height=10mm, right=8mm of r1] (h) {$H^{\otimes t}$};
  \node[hf/bus, minimum width=22mm, minimum height=26mm, right=8mm of h, yshift=-8mm] (u) {modular\\exponentiation};
  \node[hf/bus, minimum width=14mm, minimum height=10mm, right=8mm of u, yshift=8mm] (qft) {QFT};
  \node[hf/reader, minimum width=16mm, minimum height=10mm, right=8mm of qft] (m) {measure};
  \draw[hf/ask] (r1) -- (h);
  \draw[hf/ask] (h) -- (h -| u.west);
  \draw[hf/ask] (h -| u.east) -- (qft);
  \draw[hf/ask] (qft) -- (m);
  \draw[hf/answer] (r2) -- (r2 -| u.west);
  \node[hf/note, below=1mm of m, text width=26mm] {an integer close to a multiple of $q/r$};
\end{tikzpicture}
\caption{The structure of Shor's algorithm. The first register is put in
superposition, the second receives $f(a)$ for every $a$ at once, and the Fourier
transform turns the period hidden in the first register into something a
measurement can see.}
\end{figure}

\textbf{Register 1} holds the exponent $a$. It has $t$ qubits, so it can represent
$q = 2^t$ values, from $0$ to $q - 1$. Shor takes $t \approx 2\log_2 N$, for a
reason that comes at the end. Hadamards on every qubit put it in an even
superposition, while register 2 starts at zero:
\[
  \frac{1}{\sqrt q}\sum_{a=0}^{q-1} \ket{a}\ket{0}.
\]
The factor $1/\sqrt q$ keeps the norm at 1.

\textbf{Register 2} receives $f(a)$. A circuit computes $x^a \bmod N$ for the value
in register 1 and writes it into register 2. Since register 1 holds every value
at once, the result is a table of all the pairs:
\[
  \frac{1}{\sqrt q}\sum_{a=0}^{q-1} \ket{a}\,\ket{x^a \bmod N}.
\]
The period is now encoded in this state. Every $a$ that gives the same value of
$f$ is tied to it, and those $a$ are spaced exactly $r$ apart.

\section{How modular exponentiation is computed}

There is nothing mysterious in the circuit that computes $x^a \bmod N$. It is the
classical method, made reversible. Write $a$ in binary,
$a = a_0 + 2a_1 + 4a_2 + 8a_3 + \cdots$, where each $a_i$ is 0 or 1. Then
\[
  x^a = x^{a_0}\cdot x^{2a_1}\cdot x^{4a_2}\cdot x^{8a_3}\cdots,
\]
and each factor depends on a single bit: it is 1 if the bit is off and
$x^{2^i}$ if it is on. For $a = 13 = 1101_2 = 1 + 4 + 8$, this gives
$x^{13} = x\cdot x^4\cdot x^8$.

The powers $x, x^2, x^4, x^8, \dots \bmod N$ are computed in advance, classically,
by squaring repeatedly. The circuit then reads the bits of register 1 one at a
time, and each bit controls whether register 2 is multiplied by its precomputed
constant. With $t$ bits there are $t$ controlled multiplications, each on numbers
of about $\log_2 N$ bits. That is where the $(\log N)^2$ of Shor's cost comes from
(Chapter~\ref{ch:factoring}).

\begin{history}{1961--1973 \textperiodcentered{} IBM}
In 1961 Rolf Landauer showed that erasing a bit of information has an unavoidable
energy cost. Charles Bennett, also at IBM, showed in 1973 that any computation can
be made reversible, keeping enough information to run it backwards. Both results
were about the thermodynamics of ordinary computers. They became essential to
quantum computing, where every step must be unitary and therefore reversible: the
quantum circuit for $x^a \bmod N$ is Bennett's construction applied to the
schoolbook algorithm.
\end{history}

\section{Reading the period}

The last step is a quantum Fourier transform on register 1 (Chapter~\ref{ch:fourier}).
It turns a periodic pattern into peaks: the probability concentrates on values of
$k$ close to multiples of $q/r$. For $N = 15$, $x = 2$ and a register of four
qubits ($q = 16$), the period is $r = 4$, and the measurement returns $0$, $4$, $8$
or $12$, each with probability exactly $1/4$. Dividing by $q$ gives
$0, 1/4, 1/2, 3/4$: fractions with denominator $r = 4$.

When $r$ does not divide $q$, the peaks are not exact. The measured $k/q$ is then
only close to some $j/r$, and the classical method of continued fractions finds the
fraction with the smallest denominator nearby. Taking $q \geq N^2$, that is,
$t \approx 2\log_2 N$ qubits, makes the approximation good enough for continued
fractions to find $r$ with high probability. That is the reason for the size of
register 1.

\section{What the encoding decided}

Shor's choice of encoding decided everything that follows: which qubits carry which
part of the problem, how the information is processed, and how it is read out. It
also fixed the cost. Register 1 needs $O(\log N)$ qubits, modular exponentiation
needs $O((\log N)^3)$ gates in the schoolbook version, and all of it must run under
error correction, because noise anywhere in the exponentiation or the transform
corrupts the period. That cost belongs to this encoding, not to factoring itself.
Most of the progress in resource estimates since 1994 has come from finding cheaper
ways to do the same arithmetic.

\section{Three ways to encode}

Shor's is one style. There are three broad ones, each with its own costs.

\begin{concept}{Basis encoding}
Each configuration of the problem is a basis state: $n$ binary variables, $n$
qubits, and $\ket{x_1x_2\dots x_n}$ is one assignment of values. Used by Grover's
algorithm, by Shor's, and by QAOA (Farhi, Goldstone and Gutmann, 2014) for
combinatorial optimization.

\textbf{What it fixes:} reading the answer is direct. You measure and get a string
of bits.

\textbf{What it costs:} the number of qubits grows linearly with the number of
variables, and you still need interference to make the right string likely.
\end{concept}

\begin{concept}{Amplitude encoding}
The data go into the values of the amplitudes, not into which basis state is which.
A vector of $2^n$ numbers fits in $n$ qubits. Used by the HHL algorithm for linear
systems (Harrow, Hassidim and Lloyd, 2009) and by many proposals for quantum
machine learning.

\textbf{What it fixes:} an exponential compression: $2^n$ numbers in $n$ qubits.

\textbf{What it costs:} loading $2^n$ numbers into amplitudes generally takes
$2^n$ steps, and reading them back out takes many measurements. Either cost can
erase the advantage.
\end{concept}

Amplitude encoding is where most claims of exponential speedup for machine learning
die. Scott Aaronson listed the fine print of HHL in 2015: the input must load fast,
the matrix must be well behaved, and the output is a quantum state you can only
sample. In 2018 Ewin Tang, an undergraduate at the University of Texas, showed that
the most celebrated quantum recommendation algorithm could be matched by a classical
one under the same assumptions about the input. When a quantum machine learning
paper does not say how the data get into the amplitudes, assume they do not.

\begin{concept}{Hamiltonian encoding}
The problem goes into an operator, not into the initial state: build a Hamiltonian
$\hat H$ whose lowest-energy state is the solution. Used in adiabatic quantum
computing (Farhi and colleagues, 2000) and in the variational quantum eigensolver,
VQE (Peruzzo and colleagues, 2014). Andrew Lucas showed in 2014 how to write many
NP-hard problems this way.

\textbf{What it fixes:} it is natural for optimization, and for chemistry and
physics, where the Hamiltonian is the problem.

\textbf{What it costs:} the Hamiltonian must be one the hardware can implement, and
there is no guarantee that the system reaches its lowest-energy state in reasonable
time.
\end{concept}

\section{The oracle}

In many algorithms the problem shows up as an \emph{oracle}: a unitary $U_f$ that
implements the function $f$ of the problem. Given $\ket{x}$, it marks the state
according to $f(x)$, usually by flipping its sign. This is how Deutsch--Jozsa,
Bernstein--Vazirani and Grover are stated. The answers are not in the initial
state; they are in the effect the oracle has on the superposition. Here the encoding
of the problem is the construction of $U_f$.

That is also the trick hidden in Chapter~\ref{ch:qubit}: you can flip the sign of
the right answer without knowing it, as long as you can \emph{check} an answer.
The oracle is a circuit that checks. But an oracle is an assumption, not a gift.
On real hardware someone has to build $U_f$ out of gates, and its cost has to be
counted.

\section{Three costs}

Every encoding pays three costs, and they pull against each other:

\begin{itemize}
\item \textbf{qubits}: how many the mapping needs;
\item \textbf{depth}: how many sequential operations it takes to prepare the state
or apply the oracle;
\item \textbf{readout}: how much information each measurement gives, and how many
measurements you need.
\end{itemize}

Amplitude encoding uses few qubits but deep preparation and many measurements.
Basis encoding uses more qubits and reads out directly. There is no encoding that
is best in general, only one that suits the problem and the hardware. The thesis
behind this whole book is that this step, turning a real problem into states,
operators and measurements, is the hardest part of applying a quantum computer, and
the one engineers are least prepared for.
